\documentclass{article}
\usepackage{../../style/header}

\begin{document}

\title{Finite Group Theory}
\author{Lectured by Kalyani Kansal \\
Scribed by Yu Coughlin}
\date{Autumn 2025}

\maketitle

\tableofcontents

\section{Finite groups}

\begin{proposition}
    If $A,B\unlhd G$ are normal subgroups with $A\cap B = 1$ and $G=AB$, then $G\cong A\times B$.
\end{proposition}

\begin{proposition}[Orbit-stabiliser]
    If a finite group $G$ acts on a set $\Omega$, then for all $\alpha \in \Omega$, the size of the $G$-orbit of $\alpha$ is the size of $G$ divided by the size of the stabiliser.
\end{proposition}

\begin{corollary}
    When $G$ acts on itself by conjugation, for all $x\in G$, the size of the conjugacy class $x^G$ is $\abs{G:C_G(x)}$.
\end{corollary}

\subsection{Conjugacy}

\begin{proposition}
    In $S_n$, conjugacy classes are sets of permutations with the same cyclic-shape.
    \begin{proof}
        If $c = (a_1\ a_2 \ \cdots \ a_k)\in S_n$, then for all $g\in S_n$, we have $gcg^{-1} = (g(a_1) \ g(a_2) \ \cdots \ g(a_k))$. Now write $x\in S_n$ as a product of disjoint cycles, and we get $gxg^{-1}$ has the same cycle shap as $x$.
    \end{proof}
\end{proposition}

\begin{corollary}
    If $x\in S_n$ has cycle chap $(r_1^{a_1},\ldots,r_k^{a_k})$, then $\abs{C_{S_n}(c)} = \prod r_i^{a_i}a_i!$.
\end{corollary}

\begin{proposition}
    Let $x\in A_n$, if $C_{S_n}(x)\leq A_n$, then $\abs{x^{A_n}} = \frac{1}{2}\abs{x^{S_n}}$, and $x^{S_n}$ splits as two $A_n$-conjugacy classes; otherwise, $x^{A_n}=x^{S_n}$.
    \begin{proof}
        Either $C_{S_n}(x) = C_{A_n}(x)$, or $\abs{C_{A_n}(x)} = \frac{1}{2}\abs{C_{S_n}(x)}$, plug these values into orbit stabiliser.
    \end{proof}
\end{proposition}

\begin{corollary}
    $A_5$ is simple.
    \begin{proof}
        The sizes of the conjugacy classes of $A_5$ are thus $1,20,12,12,15$, the only union of these with order dividing $60$ is $1$ and $A_5$.
    \end{proof}
\end{corollary}

By year 2 linear algebra we know every $A\in GL_n(k)$ is similar to a unique matrix in RCF, so conjugacy classes are given by these.

\begin{proposition}
    If $H\leq G$, and $\Omega$ is the set of left cosets, then the map $G\rightarrow \Sym(\Omega)$ by $g\mapsto [xH\mapsto gxH]$ is a transitive action with kernel \[
    \ker(\pi) = N := \bigcap_{g\in G}gHg^{-1}\unlhd G
    \]\begin{proof}
        The first part is immediate from the definitions, and we have $x\in \ker(\pi)$ iff $g^{-1}xg\in H$ for all $g\in G$
    \end{proof}
\end{proposition}

\begin{corollary}
    If $\abs{G:N}=n$, then $\abs{G/N}\leq \im\pi\leq \Sym(\Omega)$ has order $n!$
\end{corollary}

When we set $H=1$, this recovers Cayley's theorem.

\begin{definition}
    If $H\leq G$, then $N_G(H) = \{g\in G \mid gHg^{-1} = H\}$ is called the \textbf{normaliser} of $H$ in $G$.
\end{definition}

The number of distinct conjugates of $H\leq G$ is $\abs{G:N_G(H)}$.
\begin{proposition}
    If $H\leq G$, then $N_G(H)\cong C_G(H)$ is isomorphic to a subgroup of $\Aut(H)$.\begin{proof}
        First isomorphism theorem on $\pi:N_G(H)\rightarrow \Aut(H)$ by $n\mapsto [h\mapsto nhn^{-1}]$.
    \end{proof}
\end{proposition}

\begin{proposition}
    If $H,k\leq G$ finite, then $\abs{HK} = \abs{H}\abs{K}/\abs{H\cap K}$ and if $h\leq N_G(K)$, then $HK\leq G$.
    \begin{proof}
        For the first part we count the fibres of the multiplication map $H\times K\rightarrow HK$. The second part is immediate from the definitions.
    \end{proof}
\end{proposition}

\subsection{Alternating groups}

\begin{proposition}
    For $n\geq 3$, every element of $A_n$ can be written as aproduct of $3$-cycles.\begin{proof}
        Write $x\in A_n$ as a product of an even number of $2$-cycles, $x=t_1t_2\cdots t_{2k}$, if $t_1$ and $t_2$ aren't disjoint, then their product is a $3$-cycle, otherwise, up to relabelling we can write $t_1t_2=(12)(34) = (123)(234)$.
    \end{proof}
\end{proposition}

\begin{theorem}
    For $n\geq 5$, the alternating group $A_n$ is simple.\begin{proof}
        Do induction on $n$, we already know the base case $n=5$. For $i=1,\ldots,n$ the stabiliser $(A_n)_i\cong A_{n-1}$, so if we have $N\unlhd A_n$, then as $N\cap (A_n)_i\unlhd A_{n-1}$ which, by induction, is simple, either $N\cap (A_n)_i=1$ or $(A_n)_i$.

        If $N\subset (A_n)_i$, then $N$ contains a $3$-cycle, and as the centraliser contains a disjoint transposition in $S_n\setminus A_n$, the $A_n$ orbit or a $3$-cycle equals the $S_n$ orbit which is the whole cycle-shape, so $N$ contains all $3$-cycles, and thus is all of $A_n$.

        If $N\cap (A_n)_i=1$ for all $i$, then for all $1\neq n \in N$, by relabelling we can write either $n=(12)\cdots$ or $n=(123\cdots)\cdots$, both of which don't commute with $t:=(124)$, so $[t,n]=(tnt^{-1})n^{-1}\in N$ is nontrivial, but this is $t(nt^{-1}n)=(123)(2jk)$ which moves at most $5$ elements so is in some stabiliser, a contradiction, so $N=1$.
    \end{proof}
\end{theorem}

\subsection{Matrix groups}

\begin{proposition}
    If we write $GL_n(q):=GL_n(\bF_q)$, for $q=p^m$ for some prime $p$ and $m\geq 1$, then $\abs{GL_n(q)} = (q^n-1)(q^n-q)\cdots (q^n-q^{n-1})$.
    \begin{proof}
        A matrix in $GL_n(q)$ is invertible iff all of its rows are linearly-independent in $\bF_q^n$. If we call the rows $r_1,\ldots,r_n$: there are $\abs{\bF_q^n\setminus 0} = q^n-1$ choices for $r_1$, $\abs{\bF_q^n\setminus \Span(r_1)}=q^n-q$ choices for $r_1$, $\abs{\bF_q^n\setminus\Span(r_1,r_2)} = q^n-q^2$ choices for $r_3$, and so on, until there are $\abs{\bF_q^n\setminus\Span(r_1,\ldots,r_{n-1})} = q^n-q^{n-1}$ choices for $r_n$.
    \end{proof}
\end{proposition}

\begin{proposition}
    $SL_n(q)$ is a normal subgroup of $GL_n(q)$ with order $\abs{GL_n(q)}/(q-1)$.
    \begin{proof}
        $SL_n(q)$ is the kernel of $\det:GL_n(q)\rightarrow \bF_q^\times$.
    \end{proof}
\end{proposition}

\begin{definition}
    The \textbf{projective general linear group} $PGL_n(F)$ and \textbf{projective special linear group} $PSL_n(F)$ are \[
    PGL_n(F) := GL_n(F)/Z(GL_n(F)) \qquad PSL_n(F) := SL_n(F) / Z(SL_n(F))
    \]
\end{definition}

It was a problem on the problem sheet to show these centres are \[Z(GL_n(F)) = \{\lambda I_n \mid \lambda \in F^\times\} \qquad Z(SL_n(F)) = \{\lambda I_n \mid \lambda\in \mu_n(F)\}\]

\begin{proposition}
    If $A=GL_n(F)$ with $\det(A)=\mu$, then there existst $U$ a product of elementary matrices, such that $A = U\diag(1,\ldots,\mu)$.\begin{proof}
        This is basically just row-reduction from Year 1 linear algebra. If $A=(a_{ij})$, then by adding something to row $2$ we have $a_{21}\neq 0$, and now $a_{21}^{-1}(1-a_{11}r_2)$ to $r_1$ to get $a_{11}=1$, do row reduction to the first row, and repeat by induction, this procedure never changes the determinant as elementary matrices have determinant $1$.
    \end{proof}
\end{proposition}

\begin{corollary}
    For any field $F$ and $n\geq 2$,  $SL_n(F)$ is generated by elementary matrices.
\end{corollary}

\begin{theorem}
    For $q > 3$, the group $PSL_2(q)$ is simple.
    \begin{proof}[First proof]
        Write $G=SL_2(q)$ with $q>3$, if $Z_{SL}\leq N \unlhd G$, we show $N=G$. We do this in three parts.

        Firstly, if $E_{12}(\lambda)\in N$ for some $\lambda \in \bF_q^\times$, then $N=G$. Write $U_{12} = \{E_{12}(\beta)\mid \beta \in \bF_q\}\cong(\bF_q,+)$. For any $\alpha \in \bF_q^\times$, we have \[
        \begin{pmatrix}
            \alpha & 0 \\ 0 & \alpha^{-1}
        \end{pmatrix}\begin{pmatrix}
            1 & \lambda \\ 0 & 1
        \end{pmatrix}\begin{pmatrix}
            \alpha^{-1} & 0 \\ 0 & \alpha
        \end{pmatrix} = E_{12}(\lambda\alpha^2)\in N
        \] and hence $\{E_{12}(\lambda\alpha^2) \mid \alpha\in \bF_q\}\subseteq U_{12}\cap N$. The perfect squares in $\bF_q^\times$ form a subgroup of size $(q-1)/\gcd(2,q-1)$. Ff we include the identity $E_{12}(0)$, then $\abs{U_{12}\cap N} > \frac{1}{2}q$, and hence $U_{12}\leq N$. $N$ will also contain \[
        \begin{pmatrix}
            0 & 1 \\ -1 & 0
        \end{pmatrix}\begin{pmatrix}
            1 & \lambda \\ 0 & 1
        \end{pmatrix}\begin{pmatrix}
            0 & -1 \\ 1 & 0
        \end{pmatrix} = E_{21}(\lambda)
        \] and hence $N$ contains all elementary matrices, so $N=SL_2(q)=G$.

        Secondly, if $N$ contains some \[
        A = \begin{pmatrix}
            \alpha & 0 \\ \beta & \alpha^{-1}
        \end{pmatrix}
        \] for $\alpha\in \bF_q^\times\setminus\{\pm1\}$, then $N=G$. Indeed, as $N\unlhd G$, $N$ contains the commutator $[E_{21}(1),A] = E_{21}(1-\alpha^{-2})$. So, by the previous step, $N=G$.

        Finally, we can assume $N$ contains no element of the form $A$. Let $n\in N\setminus Z_{SL}$. We assume $n$ is not triangularisable, so its minimal polynomial is an irreducible $x^2-\beta x + 1$ for some $\beta \in \bF_q$. By RCF, there exists some $y\in GL_2(q)$ such that \[
        y^{-1}ny = \begin{pmatrix}
            0 & -1 \\ 1 & \beta
        \end{pmatrix}
        \] and thus if we let $\mu = \det(y)$, and $y=g\diag(1,\mu)$ for $g\in SL_2(q)$, we get \[
        g^{-1}ng = \begin{pmatrix}
            1 & 0 \\ 0 & \mu
        \end{pmatrix}\begin{pmatrix}
            0 & -1 \\ 1 & \beta
        \end{pmatrix}\begin{pmatrix}
            1 & 0 \\ 0 & \mu^{-1}
        \end{pmatrix} = \begin{pmatrix}
            0 & -\mu \\ \mu^{-1} & \beta
        \end{pmatrix} = n'\in N
        \] and the commutator \[
        \br{
            \begin{pmatrix} 
                \alpha^{-1} & 0 \\ 0 & \alpha 
            \end{pmatrix},-n'
        } = \begin{pmatrix}
            \alpha^{-2} & 0 \\ \mu\beta(\alpha^2-1) & \alpha^2
        \end{pmatrix} = n''\in N
        \] This matrix is of the form $A$, unless $\alpha^4=1$. We can always find such an $\alpha$ when $q\neq 5$.

        If $q=5$, then let $\alpha = 2$. Then $(n'')^2$ is a non-identity elementary matrix unless $\beta = 0$, in which case $n$ is diagonalisable, a contradiction.
    \end{proof}
\end{theorem}

We can also give a more general proof, and one that better generalises to showing lots of other matrix groups are simple.

\subsection{Iwasawa's lemma}

\begin{definition}
    The action of a group $G$ on a set $\Omega$ is called \textbf{$2$-transitive} if for all pairs $(x,y),(z,w)\in \Omega^2$, with $x\neq y$ and $z\neq w$, there exists some $g\in G$ with $g(x,y)=(z,w)$.
\end{definition}

\begin{lemma}
    The action of $SL_2(F)$ on $\Omega = \{\Span(v) \mid  v \in F^2\setminus 0\}$ by $g\in SL_2(F)$ sending $\Span(v)$ to $\Span(g(v))$ is $2$-transitive, and has kernel $Z=\{\pm I_2\}$.
    \begin{proof}
        Given $(\Span(v_1),\Span(v_2))$ and $(\Span(w_1),\Span(w_2))$ two pairs of distinct lines, $v_1,v_2$ and $w_1,w_2$ are both bases for $F^2$. So there exists $g\in GL_2(F)$ sending $v_1\mapsto w_1$ and $v_2\mapsto w_2$. Take $h\in SL_2(F)$ by $v_1\mapsto \det(g)^{-1} w_1$ and $v_2\mapsto w_2$.
    \end{proof}
\end{lemma}

\begin{corollary}
    The induced action of $G=PSL_2(F)$ on $\Omega$ is $2$-transitive.
\end{corollary}

\begin{theorem}[Iwasawa's lemma]
    Let $G\leq \Sym(X)$ be a finite 2-transitive permutation group. If $G=G'$, and for all $\alpha \in X$ there exists some abelian $U\unlhd G_\alpha$ such that $\abr{gUg^{-1}\mid g\in G}=G$, then $G$ is simple.
    \begin{proof}
        Suppose we have $1\neq N \unlhd G$. Take $1\neq n \in N$, and $\alpha,\beta \in X$ such that $n(\alpha)=\beta\neq \alpha$. By 2-transitivity, for all $\gamma \in X\setminus \alpha$ there exists some $g\in G$ with $g(\alpha,\beta)=(\alpha,\gamma)$. Then $gng^{-1}\in N$ sends $\alpha$ to $\gamma$, thus $N$ is transitive.

        As $N$ is transitive, for any $\alpha \in X$, we have $NG_\alpha =G$ and for all $g\in G$, there exists $n\in N$ with $n(\alpha)=g(\alpha)$, so $n^{-1}g \in G_\alpha$, and thus $g\in NG_\alpha$.

        Let $U\unlhd G_\alpha$ be as in the hypothesis, then as $N\unlhd G$, we get $NU\unlhd NG_\alpha=G$, this thus contains $gUg^{-1}$ for all $g\in G$. By assumption, $NU$ must be all of $G$.

        We now just use the second isomorphism theorem to get $G/N = NU/N \cong U/(U\cap N)$ which is abelian, hence $G=G'\leq N$. Therefore $G=N$.
    \end{proof}
\end{theorem}

We can now give another proof that $PSL_2(q)$ is simple.

\begin{proof}[Second proof]
    Let $e_1,e_2$ be the standard basis of $\bF_q^2$, and $\alpha = \Span(e_1)\in \Omega$. All elements of $G_\alpha$ then must be of the form \[
    G_\alpha = \left\{ \begin{pmatrix}
    a & b \\ 0 & a^{-1}
    \end{pmatrix} \ \middle| \ a\in \bF_q^\times, \ b\in \bF_q\right\}/Z_{SL}
    \] So if we take \[
    U = \left\{ \begin{pmatrix}
        1 & b \\ 0 & 1
    \end{pmatrix}Z \ \middle| \ b\in \bF_q \right\} \cong (\bF_q,+)
    \] then $U$ is abelian and $U\unlhd G_\alpha$. We can also take \[
    \begin{pmatrix}
        0 & 1 \\ -1 & 0
    \end{pmatrix}\begin{pmatrix}
        1 & b \\ 0 & 1
    \end{pmatrix}\begin{pmatrix}
        0 & -1 \\ 1 & 0
    \end{pmatrix} = \begin{pmatrix}
        1 & 0 \\ b & 1
    \end{pmatrix}\in U
    \] so $\abs{gUg^{-1}}$ contains all symmetric matrices and is thus all of $PSL_2(q)$.

    Consider the commutator \[
    \begin{pmatrix}
        a & 0 \\ 0 & a^{-1}
    \end{pmatrix}\begin{pmatrix}
        1 & b \\ 1 & 0
    \end{pmatrix}\begin{pmatrix}
        a^{-1} & 0 \\ 0 & a
    \end{pmatrix}\begin{pmatrix}
        1 & -b \\ 0 & 1
    \end{pmatrix} = \begin{pmatrix}
        1 & b(a^2-1) \\ 0 & 1
    \end{pmatrix}
    \] as $q>3$ we can always find $a \in \bF_q\setminus \{0,\pm1 \}$, so $G'$ contains all of $U$. As $G'\unlhd G$, it also contains $\abr{gUg^{-1}}=G$. Hence, by Iwasawa's lemma, $G$ is simple.
\end{proof}

\subsection{Sylow's theorems}

\begin{theorem}[Sylow 1]
    If $\abs{G}=p^am$ for $p$ prime and $p\nmid m$, then $G$ has a subgroup of order $p^a$.
    \begin{proof}
        We do induction on $\abs{G}$. The base case $\abs{G}=1$ is trivial. First, for $x_1,\ldots,x_k$ representatives of the distinct non-central conjugacy classes of $G$, we have the conjugacy class equation  \[
        \abs{G} = \abs{Z(G)} + \sum_{i=1}^k \abs{x_i^G}
        \] If $p\mid \abs{Z(G)}$, then Cauchy's theorem gives some $z\in Z(G)$ of order $p$. As $\abr{z}\unlhd G$, we can find $H\leq G/\abr{z}$ of order $p^{a-1}$. And $H=P/\abr{z}$ for some $P\leq G$, so $\abs{P}=p^{a-1}$.

        If $p\nmid \abs{Z(G)}$, then one of the $x_i^G$ has $p\nmid \abs{x_i^G}$. As $\abs{x_i^G}=\abs{G:C_G(x_i)}$, we must have $p^a\mid \abs{C_G(x_i)}$. Thus, by induction, there exists $P\leq C_G(x_i) < G$ with $\abs{P}=p^a$.
    \end{proof}
\end{theorem}

\begin{definition}
    Such a subgroup is called a \textbf{Sylow $p$-subgroup}. Write $\Syl_p(G)$ for the set of all Sylow $p$-subgroups. $G$ acts by conjugation on $\Syl_p(G)$, the stabiliser of $P\in \Syl_p(G)$ is $N_G(P)$. Wrtie $n_p(G) := \abs{\Syl_p(G)}$.
\end{definition}

\begin{lemma}
    If $P\in \Syl_p(G)$, $Q$ is a $p$-subgroup of $G$, and $Q\leq N_G(P)$, then $Q\leq P$.
    \begin{proof}
        If $Q\leq N_G(P)$, then $PQ\leq G$, and $\abs{PQ}=\abs{P}\abs{Q}/\abs{P\cap Q} = p^a\abs{Q:P\cap Q}$, as $Q$ is a $p$-group, and $p^{a+1}\nmid \abs{G}$, we have $Q=P\cap Q$, thus $Q\leq P$.
    \end{proof}
\end{lemma}

\begin{theorem}[Sylow 2]
    For all primes $p$, we have $n_p(G) \equiv 1$ (mod $p$).
    \begin{proof}
        For $P\in \Syl_p(G)$, consider the action of $P$ on $\Syl_p(G)$ by conjugation. If $\{Q\}$ is any singleton orbit, then $P\leq N_G(Q)$, so $P\leq Q$ and hence $P=Q$. So this action has exatly one singleton orbit. Now, by the orbit stabiliser theorem, the other orbits have size dividing $\abs{P}$, so are divisible by $p$. As $n_p(G)=\abs{Syl_p(G)}$ is a sum of orbits, $n_p(G)\equiv 1$ (mod $p$).
    \end{proof}
\end{theorem}

\begin{theorem}[Sylow 3]
    If $Q$ is a $p$-subgroup of $G$, then there exists $P\in \Syl_p(G)$ with $Q\leq P$.
    \begin{proof}
        Consider the action of $Q$ on $\Syl_p(G)$ by conjugation. As $n_p(G)\equiv 1$ (mod $p$), not every $Q$-orbit has size divisible by $p$, so there is an orbit of size $1$, call it $\{P\}$. Then $Q\leq N_G(P)$, so $Q\leq P$.
    \end{proof}
\end{theorem}

\begin{theorem}[Sylow 4]
    The action of $G$ on $\Syl_p(G)$ by conjugation is transitive.
    \begin{proof}
        Let $P\in \Syl_p(G)$, and $\Omega$ be the $G$-oribt of $P$ in $\Syl_p(G)$. Again, there is some orbit of size $1$, thus $\abs{\Omega}\equiv 1$ (mod $p$). If $Q\in \Syl_p(G)$, consider the action of $Q$ on $\Omega$. Again, there is some orbit of size $1$, call it $\{R\}$. So $Q\leq N_G(R)$, hence $Q =R\in \Omega$.
    \end{proof}
\end{theorem}

\begin{corollary}
    If $P\in \Syl_p(G)$, then $n_p(G) = \abs{G:N_G(P)}$. In particular, $n_p(G)\mid \abs{G}$ and $n_p(G)=1$ iff $P\unlhd G$.
    \begin{proof}
        First, the stabiliser of $P$ in $\Syl_p(G)$ is $N_G(P)$, so this is just orbit stabiliser. The others follow.
    \end{proof}
\end{corollary}

\begin{proposition}
    If $\abs{G} = pq$ for $p>q$ primes, then $G$ has a normal Sylow $p$-subgroup. If, furthermore, $q\nmid p-1$, then $G\cong C_{pq}$.
    \begin{proof}
        First, $n_p(G)\leq \abs{G}/p=q$, and $n_p(G)\equiv 1$ (mod $p$). As $q<p$, we thus have $n_p(G)=1$. If $q\nmid p-1$, then $p\not\equiv 1$ (mod $q$). So, similarly, we get a normal Sylow $q$-subgroup. Call these $P,Q\unlhd G$. As $P\cap Q = 1$ and $G=PQ$, we must have $G\cong P\times Q \cong C_{pq}$.
    \end{proof}
\end{proposition}

\begin{proposition}
    If $\abs{G}=p^2q$ with $p,q$ primes, then $G$ has either a normal Sylow $p$-subgroup or a normal Sylow $q$-subgroup.
    \begin{proof}
        If $n_p(G),n_q(G)\neq 1$, then we can't have $n_q(G)=p$, as $p\not\equiv 1$ (mod $q$), by $n_p(G)=q\equiv 1$ (mod $p$). So $n_q(G)=p^2$. These must be $p^2$ disjoint copies of $C_q$, so $\abs{G}-p^2$ elements of order $q$ in $G$. This only leaves room for one Sylow $p$-subgroup.
    \end{proof}
\end{proposition}

\begin{proposition}
    Let $\abs{G}=p^a m$, where $p$ is prime, $a\geq 1$, $m\geq 2$, and $p\nmid m$. If $G$ is simple, then $\abs{G}\mid n_p(G)!$.\begin{proof}
        As $G$ acts transitively on $\Syl_p(G)$, there is a homomorphism $\phi:G\rightarrow \Sym(\Syl_p(G))$. The kernel is a proper normal subgroup of $G$, thus $\phi$ is injective. So $G$ is isomorphism so a subgroup of $S_{n_p(G)}$.
    \end{proof}
\end{proposition}

\subsection{Extenstions and semidirect products}

\begin{definition}
    An extension of $N$ by $H$ is a group $G$ with $N\cong N_0 \unlhd G$ such that $G/N_0\cong H$.
\end{definition}

The Jordan-Holder theorem says that finding all extensions of $N$ by $H$ is the second half of classifying all finite groups. After we know the finite simple groups.

\begin{definition}
    If $G$ is an extensions of $N$ by $H$, we say $G$ is \textbf{split} if there is some $H_0\leq G$ such that $G=N_0H_0$ and $N_0\cap H_0=1$. Note that $H=G/N_0 = N_0H_0/H_0\cong H_0$.
\end{definition}

\begin{proposition}
    if $G=NH$ with $N\cap H=1$ is a split extension, then we have the following: \begin{enumerate}
        \item every $g\in G$ can abe written as $g=nh$ for unique $n\in N$ and $h\in H$;
        \item for all $h\in H$ the map $\iota_h:N\rightarrow N$ by $\iota_h(n)=hnh^{-1}$ is in $\Aut(N)$;
        \item the map $\iota:H\rightarrow \Aut(N)$ sending $h\mapsto \iota_h$ is a homomorphism;
        \item multiplication is determined by $\iota$ as $(n_1h_1)(n_2h_2)=(n_1\iota_{h_1}(n_2))(h_1h_2)$.
    \end{enumerate}
    \begin{proof}
        These are all immediate.
    \end{proof}
\end{proposition}

\begin{definition}
    Given $\iota:H\rightarrow \Aut(N)$, we give the set $N\times H$ the multiplication $(n_1,h_1)(n_2,h_2) = (n_1\iota_{h_1}(n_2),h_1h_2)$, where we write $\iota_{h_1}:=\iota(h_1)$. Call this the \textbf{semidirect product} $N \rtimes_\iota H$.
\end{definition}

\begin{proposition}
    $N\rtimes_\iota H$ is a group under this multiplication, and a split extension of $N\times\{1\}$ by $\{1\}\times H$. If $\iota$ is the map from the previous proposition, then $G\cong N \rtimes_\iota H$.
    \begin{proof}
        These are also all immediate.
    \end{proof}
\end{proposition}

Given to group $N,H$, if we find all homomorphisms $\iota:H\rightarrow \Aut(N)$, and compute all the corresponding semidirect products $N\rtimes_\iota H$, then these are all the split extensions, up to isomorphism.

\begin{proposition}
    Let $N,H$ be finite groups, with $\iota,\kappa:H\rightarrow \Aut(N)$ homomorphism. If there exists $\alpha \in \Aut(N)$ and $\beta \in \Aut(H)$ such that $\kappa(\beta(h))=\alpha\iota(h)\alpha^{-1}$ for all $h\in H$, then $N\rtimes_\iota H \cong N\rtimes_\kappa H$.\begin{proof}
        Consider $\phi:N\rtimes_\iota H \rightarrow N \rtimes_\kappa H$ by $\phi(n,h)=(\alpha(n),\beta(h))$. This is an isomorphism.
    \end{proof}
\end{proposition}

This condition is sufficient but certainly not necessary.

\begin{proposition}
    If $N$ and $H$ are finite abelian groups of coprime order, then every extension of $N$ by $H$ splits.
    \begin{proof}
        Exercise.
    \end{proof}
\end{proposition}

\section{Series}

\begin{definition}
    A subgroup $N\leq G$ is \textbf{characteristic} if for all $\alpha \in \Aut(G)$, $\alpha(N)=N$, written $N\char G$. These are clearly automatically normal.
\end{definition}

\begin{proposition}
    If $M\char N \unlhd G$ are subgroups, then $M\unlhd G$.
    \begin{proof}
        For all $g\in G$, conjugation by $g$ is an automorphism of $N$, under which $M$ is invariant.
    \end{proof}
\end{proposition}

The centre, commutator, and inner automorphism group are all characteristic.

\subsection{Composition series}

\begin{definition}
    A proper normal subgroup $N\lhd G$ is called \textbf{maximal} if there are no proper $N\subsetneq M\lhd G$.
\end{definition}

\begin{lemma}
    $N$ is a maximal normal subgroup of $G$ iff $G/N$ is simple.\begin{proof}
        This is by the bijective correspondence between normal subgroups of $G/N$ and normal subgroups of $G$ containing $N$.
    \end{proof}
\end{lemma}

\begin{definition}
    A \textbf{composition series} for a group $G$ is a sequence of maximal subgroups \[
    G = G_0 \rhd  G_1 \rhd \cdots \rhd G_{r-1} \rhd G_r = 1
    \] where the simple quotient $G_i/G_{i-1}$ are called the \textbf{composition factors}.
\end{definition}

\begin{proposition}[Second isomorphism theorem]
    If $H\leq G$ and $N\unlhd G$, then $H\cap N \unlhd H$, $HN\leq G$, and $HN/N\cong H/(H\cap N)$.
    \begin{proof}
        The first two parts are clear, for the third consider the map $\phi:H\rightarrow HN/N$ by $\phi(h) = hN$, this is a surjective homomorphism and the kernel can be checked to be exactly $H\cap N$, the result now comes from the first isomorphism theorem.
    \end{proof}
\end{proposition}

\begin{theorem}[Jordan-H\"older]
    The composition factors of a group are well-defined up to reordering.
    \begin{proof}
        Do induction on $\abs{G}$, the base case is trivial. Suppose $G$ has two composition series: \[
        G \rhd G_1 \rhd G_2 \rhd \cdots \rhd G_r = 1 \qquad G \rhd H_1 \rhd H_2 \rhd \cdots \rhd H_s = 1
        \] then if $G_1=H_1$, the two remaining series are composition series for a strictly smaller group, so have the same compositions factors, by adding $G/G_1$ we are done.

        Suppose, instead, that $G_1\ncong H_1$, then let $K = G_1\cap H_1\lhd G$ and choose a composition series \[
        K \rhd K_3 \rhd \cdots \rhd K_t = 1
        \] I now claim $G_1/K\cong G/H_1$ and $H_1/K\cong G/G_1$, we have $G_1\subset G_1H_1\unlhd G$, so $G_1H_1 = G$ and by the second isomorphism theorem $G/H_1 = G_1H_1/H_1\cong G_1 / (G_1\cap H_1) = G_1/K$. As $G_1/K$ is simple, we have the two composition series for $G_1$ \[
        G_1 \rhd G_2 \rhd G_3 \rhd \cdots \rhd G_r = 1 \qquad G_1 \rhd K \rhd K_3 \rhd \cdots \rhd K_t = 1
        \] which, by induction, must have the same length and composition factors, if we do the same for $H_1$, then we can rewrite our initial composition factors as \begin{gather*}
            \{G/G_1,G_1/G_2,G_2/G_3,\ldots,G_{r-1}\} = \{G/G_1,G_1/K,K/K_3,\ldots, K_{t-1}\} \\ 
            = \{H_1/K,G/H_1,K/K_3\ldots, K_{t-1}\} = \{G/H_1,H_1/H_2,H_2/H_3,\ldots, H_{s-1}\}\qedhere
        \end{gather*}
    \end{proof}
\end{theorem}

\newpage

\subsection{Solvable groups}

\begin{definition}
    A finite group $G$ is \textbf{solvable} if it has a series of subgroups \[
    1 = N_0 \leq N_1 \leq \cdots \leq N_s = G
    \] where $N_i\unlhd G$ and $N_i/N_{i-1}$ is abelian for all $i$.
\end{definition}

Abelian groups are solvable. Finite $p$-groups are also solvable, as they always have nontrivial centre.

\begin{definition}
    The \textbf{$n$th derived subgroup} $G^{(n)}$ of $G$ is defined recursively as $G^{(0)}:= G$ and $G^{(i+1)} := (G^{(i)})'$. The \textbf{derived series} of $G$ is then the series  $G\geq G^{(1)} \geq G^{(2)} \geq \cdots$.
\end{definition}

\begin{proposition}
    If $\phi:G\rightarrow H$ is a surjective homomorphism, then $\phi(G^{(i)}) = H^{(i)}$.\begin{proof}
        For $x,y\in G$, we have $\phi([x,y]) = [\phi(x),\phi(y)]$, so $\phi$ being surjective is enough.
    \end{proof}
\end{proposition}

\begin{corollary}
    For all $i$, we have $G^{(i)}\char G$.
\end{corollary}

\begin{proposition}
    $G$ is solvable iff $G^{(r)} = 1$ for some $r$.
    \begin{proof}
        If $G^{(r)}=1$, the consider the derived series \[
        1 = G^{(r)} \leq G^{(r-1)} \leq \cdots \leq G^{(1)} \leq G
        \] where each $G^{(i)}\unlhd G$ and each quotient $G^{(i)}/G^{(i+1)} = G^{(i)}/(G^{(i)})'$ is abelian. So $G$ is solvable.

        Conversely, if $G$ is solvable, with series \[
        1 = N_0 \leq N_1 \leq \cdots \leq N_s = G
        \] where $N_i\unlhd G$ and $N_i/N_{i-1}$ is abelian for all $i$, then $N_i'\leq N_{i-1}$, so $G'\leq N_{s-1}$, $G^{(2)} \leq N_{s-1}' \leq N_{s-2}$ and so on. So, eventually, $G^{(s)} \leq N_0 = 1$.
    \end{proof}
\end{proposition}

\begin{definition}
    For a solvable group $G$, the \textbf{derived length} of $G$ is $\dl(G):=\min\{r \mid G^{(r)}=1\}$.
\end{definition}

\begin{lemma}
    Subgroups and quotient groups of solvable groups are solvable.
    \begin{proof}
        If $G$ is solvable, and $H\leq G$, then $H'\leq G'$, and thus $H^{(i)}\leq G^{(i)}$. Hence $H^{(i)} = 1$ for some $i$, so $H$ is solvable.

        Now, if $N\unlhd G$, and we consider the surjective quotient map $\phi:G\rightarrow G/N$, then $(G/N)^{(i)} = \phi(G^{(i)}) = 1$ for some $i$, hence $G/N$ is solvable.
    \end{proof}
\end{lemma}

\begin{lemma}
    If $N\lhd G$ and both $N$ and $G/N$ are solvable, then $G$ is solvable.
    \begin{proof}
        
    \end{proof}
\end{lemma}

\begin{proposition}
    If $G$ is a finite group with composition series \[
    G = G_0 \rhd  G_1 \rhd \cdots \rhd G_{r-1} \rhd G_r = 1
    \] then $G$ is solvable iff all the composition factors $G_i/G_{i-1}$ are cyclic and of prime order.
    \begin{proof}
        
    \end{proof}
\end{proposition}

\begin{definition}
    If $A$ is a finite abelian group and there exists a prime $p$ such that for all $x\in A$, $x^p=1$, then we call $A$ \textbf{elementary abelian}.
\end{definition}

By the classification of finite abelian groups, the elementary abelian groups are $\bF_p^n$.

\begin{definition}
    A normal subgroup $N\unlhd G$ is called \textbf{minimal normal} if $N$ contains no smaller nontrivial normal subgroup of $G$.
\end{definition}

\begin{proposition}
    Let $G$ is a finite group, with $N$ a minimal normal subgroup. If $N$ is solvable, then $N$ is also elementary abelian.
\end{proposition}

\begin{definition}
    For a finite group $G$, a \textbf{cheif series} is a series \[
    1 = N_0 \leq N_1 \leq \cdots \leq N_r = G
    \] where each $N_i\unlhd G$, and each $N_{i+1}/N_i$ is minimal normal in $G/N_i$. The quotient $N_{i+1}/N_i$ are called \textbf{cheif factors}.
\end{definition}

While there is a Jordan-H\"older type theory for the uniqueness of the set of chief factors of a finite group, we won't prove it here.

\begin{corollary}
    The chief factors of a finite solvable group are all elementary abelian.
\end{corollary}

\subsection{Hall's theorem}

\begin{definition}
    Let $\pi$ be a set of primes, we say a positive integer $n$ is a \textbf{$\pi$-number} if all the prime factors of $n$ are in $\pi$. In this context, write $\pi'$ for the complement of $\pi$. A subgroup $H\leq G$ is called a \textbf{$\pi$-subgroup} if $\abs{H}$ is a $\pi$-number.
\end{definition}

\begin{proposition}[Frattini Argument]
    Let $N\unlhd G$ and $P\in\Syl_p(N)$, then $G=N_G(P)N$.
    \begin{proof}
        For all $g\in G$, we have ${}^gP\leq{}^gN$, thus ${}^gP\in\Syl_p(N)$ and by Sylow $4$, there exists some $n\in N$ with ${}^gP={}^nP$, thus $gn^{-1}\in N_G(P)$, and $g\in N_G(P)N$.
    \end{proof}
\end{proposition}

The following is a generalisation of a question on an earlier problem sheet. These are both special cases of the Schur-Zassenhaus theorem, which applies to arbitrary groups, and will not be proved here.

\begin{theorem}
    If $N$ and $H$ be finite solvable groups of coprime orders, then every extension of $N$ by $H$ splits.
\end{theorem}

\begin{theorem}[Hall 1]
    If $\pi$ be any set of primes, and $G$ a finite solvable group with $\abs{G}=nm$ where $n$ is a $\pi$-number and $m$ a $\pi'$-number, then $G$ has a subgroup of order $n$.
\end{theorem}

\begin{definition}
    Such a subgroup $H$ is called a \textbf{Hall $\pi$-subgroup} of $G$.
\end{definition}

Like we have generalised Sylow 1 to Hall subgroups, there are similar generalisations of Sylow 3 and Sylow 4 to Hall subgroups, we won't prove these here.

\newpage

\subsection{Nilpotent groups}

\begin{definition}
    A finite group $G$ is \textbf{nilpotent} if there exists a series of subgroups: \[
    1 = G_0 \leq G_1 \leq \cdots \leq G_r = G\]
    such that for all $0\leq i \leq r-1$ \[
    G_i\unlhd G \text{ and } G_{i+1}/G_i \leq Z(G/G_i)
    \] such a series of subgroups is called a \textbf{central series}
\end{definition}

Notably, all abelian groups are nilpotent; and all nilpotent groups are solvablem, as $G_{i+1}/G_i\leq Z(G/G_i)$ implies $G_{i+1}/G_i$ is abelian.

\begin{proposition}
    Every finite $p$-group is nilpotent.
    \begin{proof}
        Let $P$ be such a nontrivial $p$-group, define a series of subgroups as follows:\begin{itemize}
            \item $P_0:= 1$;
            \item $P_1 := Z(P)$;
            \item Have $P_2$ a subgroup containing $P_1$ such that $P_2/P_1 = Z(P/P_1)$;
            \item recursively define each $P_{i+1}$ to be a sugroup containing $P_i$ such that $P_{i+1}/P_i = Z(P/P_i)$.
        \end{itemize}
        as nontrivial $p$-groups always have nontrivial centre this series is strictly increase, so must termiante at some $P_r=P$.
    \end{proof}
\end{proposition}

\begin{definition}
    The \textbf{upper central series} of a finite group $G$ is the central series $Z_i(G)$ formed in the above proposition. We are taking the \textit{biggest step} possible at each point.
\end{definition}

\begin{definition}
    For $G$ a finite group and $H,K\subseteq G$, define \[
    [H,K] := \abr{[h,k] \mid h\in H, k\in K}\leq G
    \] one can immediately see $[H,K]=[K,H]$. We want to define the central series that takes the biggest step down at each point: this will be the \textbf{lower central series} of $G$ defined as follows:\begin{itemize}
        \item $\gamma_0(G) := G$;
        \item For $\gamma_1(G)$ to be the next step down it must satisfy \[
        G/\gamma_1(G) \leq Z(G/\gamma_1(G))
        \] which just means $G/\gamma_1(G)$ is abelian, so we must have $[G,G]\leq \gamma_1(G)$, if we want to make the biggest just down $\gamma_1(G)$ must be as small as possible so we set $\gamma_1(G):=[G,G]$;
        \item Now given our $\gamma_i(G)$, to construct $\gamma_{i+1}(G)$ we must satisfy: \[
            \gamma_i(G)/\gamma_{i+1}(G) \leq Z(G/\gamma_{i+1}(G))
        \] so $\gamma_i(G)$ should commutate with everything in $G$ up to $\gamma_{i+1}(G)$ which is just: \[
            [G,\gamma_i(G)]\leq \gamma_{i+1}(G)
        \] so to continue with our biggest steps, we should set $\gamma_{i+1}(G) := [\gamma_i(G),G]$.
    \end{itemize}
\end{definition}

We shoul expect our two intuitively shortest possible sequences to be of the same minimal length, i.e.~our group is nilpotent iff these series terminate. This is true, we will first need a lemma to prove it.

\begin{lemma}
    For $G$ a finite group and $X,Y\leq G$: \begin{enumerate}
        \item $Y\leq N_G(X)$ iff $[X,Y]\leq X$;
        \item If $Y\unlhd G$ and $Y\leq X$ then we have\[
            X/Y \leq Z(G/Y) \iff [X,G]\leq Y
        \]
    \end{enumerate}\begin{proof}
        (1) $Y\leq N_G(X)$ iff for all $x\in X$ and $y\in Y$ we have $[x,y]=xyx^{-1}y^{-1}\in X$ iff $[X,Y]\leq X$.

        (2) For all $x\in X$ we have: $xY\in Z(G/Y)$ iff $[xY,gY] = Y$ for all $g\in G$ iff $[x,g]Y=Y$ iff $[x,g]\in Y$.
    \end{proof}
\end{lemma}

This formalises our intuition of taking the biggest steps possible in the definitions of our upper and lower central series.

\begin{theorem}
    For $G$ a finite group and $n$ a positive integer, $Z_n(G)=G$ iff $\gamma_n(G)=1$. Furthermore, $G$ is nilpotent iff these conditions hold for some $n$.
    \begin{proof}
        If either of the statements are true, then $G$ is certainly nilpotent. So it suffices to prove for nilpotent $G$ the set of $n$ such that $Z_n(G) = G$ and $\gamma_n(G)=1$ are equal and nonempty.

        Let \[
        1 = N_0 \leq N_1 \leq \cdots \leq N_k = G
        \] be a central series for $G$. By the previous lemma, for each $i$ we have $[N_{i+1},G]\leq N_i$.

        Consider $\gamma_0(G)=G\leq N_k$, so $\gamma_1(G)=[G,G]\leq[N_k,G]\leq N_{k-1}$ and inductively: \[
        \gamma_i(G) = [\gamma_{i-1}(G),G] \leq [N_{k-i+1},G]\leq N_{k-1}
        \] and $N_0=1$ we must have $\gamma_k(G)=1$, applying this to our upper central series is half way.

        Now consider $Z_0(G) = 1 = N_0(G)$, assume for $i\geq 1$ that $N_{i-1}\leq Z_{i-1}(G)$, equivalently: \[
        [N_i,G]\leq N_{i-1}\leq Z_{i-1}(G)
        \] moreover $[N_iZ_{i-1}(G),G]\leq Z_{i-1}$ as for all $n\in N_i$, $z\in Z_{i-1}(G)$, and $g\in G$ we have \[
        [nz,g]Z_{i-1}(G) = [n,zg][z,g]Z_{i-1}(G)\subseteq [N_i,G]Z_{i-1}(G) \leq Z_{i-1}(G)
        \] hence \[
            N_iZ_{i-1}(G)/Z_{i-1}(G) \leq Z(G/Z_{i-1}(G)) = Z_{i}(G)/Z_{i-1}(G)
        \] so $N_i\leq Z_i(G)$ and by induction we are done.
    \end{proof}
\end{theorem}

\begin{definition}
    The a niplotent group $G$ the smallest integer $n$ such that $Z_n(G)=G$ (or equivalently $\gamma_n(G)=1$) is called the nilpotentce class of $G$.
\end{definition}

\begin{corollary}
    Subgroups and quotients of nilpotent groups are also nilpotnet.\begin{proof}
        Let $G$ have nilpotentcy class $n$, if $H\leq G$ then $\gamma_n(H)\leq\gamma_n(G)=1$ so $H$ is nilpotent.

        Let $\phi$ be the quotient map $G\rightarrow G/N$ for some $N\unlhd G$, as homomorphisms commute with commutators and $\phi$ is surjective, we have $\phi(\gamma_i(G))=\gamma_i(\phi(G))$ for all $i$, specifically $\gamma_n(\phi(G))=1$ so $G/N$ is nilpotent.
    \end{proof}
\end{corollary}

If such an $N$ and $G/N$ are both nilpotent, $G$ is not necessarily nilpotent. For example take $S_3=C_3\rtimes C_2$ which isn't nilpotent: $[S_3,S_3]=C_3$, $[C_3,S_3]=C_3$.

Nilpotent groups have lots of very nice characterisations.

\begin{theorem}
    For a finite group $G$ tfae: \begin{enumerate}
        \item $G$ is nilpotent;
        \item for any $H < G$, we have $H < N_G(H)$ is proper;
        \item every maximal subgroup of $G$ is normal;
        \item every Sylow subgroup of $G$ is normal;
        \item $G\cong P_1\times \cdots \times P_k$ for $P_i$ a $p_i$-group with all the $p_1,\ldots,p_k$ distinct primes.
    \end{enumerate}

    \begin{proof}
        ($1\Rightarrow 2$) Let $1=N_0\leq \cdots \leq N_k = G$ be a central series, and let $H < G$. Let $r$ be the largest positive integer such that $N_r\leq H$, then $N_{r+1}/N_r\leq Z(G/N_r)$ and equivalently [$N_{r+1},H]\leq N_r \leq H$, by our lemma we know this is equivalent to $N_{r+1}\leq N_G(H)$, hence $H\leq N_G(H)$.

        ($2\Rightarrow 3$) Let $M$ be a maximal subgroup of $G$, then $M\leq N_G(M)=G$, as $M$ is maximal.

        ($3\Rightarrow 4$) Let $P\in \Syl_p(G)$ with $P$ not normal in $G$, so $N_G(P) < G$. There exists a maximal subgroup $M$ of $G$ containing $N_G(P)$ with $M\lhd G$. As $P\in\Syl_p(M)$ we can use the Frattini argument to get $G=N_G(P)M = M$, a contradiction, thus $P\unlhd G$.

        ($4\Rightarrow 5$) Let $P_1,\ldots P_k$ be the Sylow subgroups of $G$ for distinct primes, as each $P_i$ is normal in $G$, $[P_i,P_j]\leq P_i\cap P_j = 1$, hence $P_1P_2\cdots P_k = G$ with each $g$ written uniquely as a product of elements from different $P_i$.

        ($5\Rightarrow 1$) Products commute with the $\gamma_i$, so the finite maximum of all the components nilpotency classes will be the nilpotency class of $G$.
    \end{proof}
\end{theorem}

\subsection{Solvable radical and fitting subgroup}

\section{Applications}

\subsection{The Frattini subgroup}

\begin{definition}
    The $G$ a finite group the intersection all all maximal subgroups is called the \textbf{Frattini subgroup} and denoted $\Phi(G)$.
\end{definition}

\begin{proposition}
    For all finite groups $G$ and $x\in G$ the following are equivalent: \begin{enumerate}
        \item $x\in \Phi(G)$;
        \item For all $S\subseteq G$, if $G=\abr{S,x}$, then $G=\abr{S}$.
    \end{enumerate}
    \begin{proof}
        ($1\Rightarrow 2$) Suppose $x\in\Phi(G)$ and there exists some $S\subset G$ such that $\abr{S,x}=G$ but $\abr{S}\neq G$, then there exists some maximal subgroup $M< G$ containing $S$, but $x\in\Phi(G)\leq M$, thus $G=\abr{S,x}\leq M$, a contradiction.

        ($2\Rightarrow 1$) If $x\in G$ satisfies property $2$ and $x\notin\Phi(G)$ then there exists some maximal $M<G$ with $x\notin M$, thus $\abr{M,x}=G$, but $\abr{M}\neq G$, a contradiction.
    \end{proof}
\end{proposition}

We call elements of $\Phi(G)$ the \textbf{non-generators} of $G$.

\begin{proposition}
    $\Phi(G)$ is a characteristic nilpotent subgroup of $G$.
    \begin{proof}
        An automorphism of $G$ must permute the maximal subgroups, so fixes their intersection, thus $\Phi(G)\char G$.

        Let $p$ be a prime and $P\in\Syl_p(\Phi(G))$, by the Frattini argument, we must have $G=N_G(P)\Phi(G)$. So by the previous proposition (and some induction) $G=N_G(P)$ so $P\unlhd G$, and a fortiori $G\unlhd\Phi(G)$. Now, by the characterisation of nilpotent subgroups, $\Phi(G)$ is nilpotent.
    \end{proof}
\end{proposition}

\begin{theorem}[Burnside Basis Theorem]
    For $P$ a $p$-group with $\abs{P}=p^n$ and $\abs{\Phi(P)} = p^{n-d}$ for some $d\geq 1$, we have the following: \begin{enumerate}
        \item $P/\Phi(P)\cong \bF_p^d$ the elementary abelian group of order $p^d$;
        \item If $N\unlhd P$ and $P/N$ is an elementary abelian group, then $\Phi(P)\leq N$, thus $\Phi(G)$ is the unique normal subgroups with minmal elementary abelian quotient;
        \item For $x_1,\ldots,x_d\in P$ which equivalence classes $\overline{x_i}\in P/\Phi(P)$ respectively: \[
            P = \abr{x_1,\ldots,x_d} \iff \overline{x_1}, \ldots, \overline{x_d} \text{ is a basis for $\bF_p^d$}
        \]
    \end{enumerate}
    \begin{proof}
        (1) We know any maximal subgroup $M<P$ is normal, for a maximal normal subgroup the quotient is simple, so $\abs{P/M}=p$. Thus $P'\leq M$ and $x^p\in M$ for all $x\in P$, so $P'\leq\Phi(P)$ and $P/\Phi(P)$ is elementary abelian.

        (2) Suppose $N\unlhd P$ and $P/N$ is elementary abelian, then $\Phi(P/N)=1$. The maximal subgroups of $P/N$ are of the form $M/N$ for some maximal $M<P$ containing $N$. As $\Phi(P/N)=N$, the intersection of these must be contained in $N$, thus $\Phi(P)\leq N$.

        (3) If $x_1,\ldots, x_d$ span $P$, then they must span $\bF_p^d$ and so form a basis. If $\overline{x_1}\ldots,\overline{x_d}$ is a basis for $\bF_p^d$ then $P=\abr{x_1,\ldots,x_d,\Phi(P)}$ and thus $P = \abr{x_1,\ldots,x_d}$.
    \end{proof}
\end{theorem}

We can use this to relate the minimal number of generators of $P$ to its Frattini subgroup:

The minimal number of generators of $D_8$ is $2$, so $D_8/\Phi(D_8)\cong C_2^2$, the only order two subgroup $\Phi(D_8)$ can be is $\abr{r^2}$.

\subsection{Groups of order $p^3$}

The only ones left are non-abelian groups of order $p^3$ with $p$ odd.

\begin{lemma}
    Let $G$ be a group such that $G'\leq Z(G)$, then for all $x,y\in G$ and $n\geq 1$: \[
        x^ny^n = (xy)^n[x^{-1},y^{-1}]^{n(n-1)/2}
    \]\begin{proof}
        We go by induction on $n$, noting:\[
        xy = yx[x^{-1},y^{-1}]
        \] and that by assumption $[x^{-1},y^{-1}]\in Z(G)$. Assume the claim is true for $n$ and consider: \begin{align*}
            (xy)^{n+1}[x^{-1},y^{-1}]^{n(n+1)/2} &= (xy)^n[x^{-1},y^{-1}]^{n(n-1)/2}xy[x^{-1},y^{-1}]^n \\
            &= x^ny^nx[x^{-1},y^{-1}]^ny = x^ny^{n-1}xy[x^{-1},y^{-1}]^{n-1}y \\
            &\ \vdots \\
            &= x^{n+1}y^{n+1}\qedhere
        \end{align*}
    \end{proof}
\end{lemma}

\begin{theorem}
    For $p$ an odd prime, there are two non-abelian groups of order $p^3$ up to isomorphism: \begin{enumerate}
        \item $P_1 = \abr{a,b,c \mid a^p=b^p=c^p=1,\ [a,b]=c,\ [a,c]=[b,c]=1}\cong \abr{a,c}\rtimes\abr{b}\cong C_p^2\rtimes C_p$;
        \item $P_2 = \abr{a,b \mid a^{p^2}=b^p=1,\ [a,b]=a^p} \cong \abr{a}\rtimes\abr{b}\cong C_{p^2}\rtimes C_p$.
    \end{enumerate}
    \begin{proof}
        Let $P$ be such a group, then By the Burnside basis theorem $P/\Phi(P)\cong C_p^d$ for one of $d=1,2,3$. If $d=1$, then $P$ is cyclic; it $d=3$ then $P\cong C_p^3$ abelian, so we must have $d=2$. Write: \[
        P = \abr{a,b}, \quad \Phi(P) = \abr{c}
        \] for some $a,b,c\in P$. We also know $\Phi(P) = Z(P) = P'$.

        First, assume $a^p=b^p=1$. As $[a,b]\neq 1$ ($P$ non-abelian), take $c=[a,b]$, this gives all the relations for $P_1$, so this caes is done.

        Secondly, assume $a$ has order $p^2$, from the proof of the Burnside basis theorem we know $a^p\in \Phi(P)$, so $\Phi(P) = \abr{a^p}$. If $b^p=1$, then by taking $[a,b]=a^p$we have all the relations for $P_2$.

        If, instead, we let $b$ have order $p^2$, wlog we may assume $a^p = c$ and $b^p = c^{-1}$ (by choosing different powers of $a$ and $b$). Now by the previous lemma we can write \[
        1 = a^pb^p = (ab)^p[a^{-1},b^{-1}]^{p(p-1)/2}
        \] where $p$ being odd must divide $p(p-1)/2$, thus $[a^{-1},b^{-1}]^{p(p-1)/2}=1$ and hence $(ab)^p=1$ and by replacing $b$ with $ab$ we reach the first state.
    \end{proof}
\end{theorem}

$P_1$ also has an upper unitriangular representation over $\bF_p^3$ given by: \[
a \mapsto \begin{pmatrix}
    1 & 1 & 0 \\
    0 & 1 & 0 \\
    0 & 0 & 1
\end{pmatrix}, \quad
b \mapsto \begin{pmatrix}
    1 & 0 & 0 \\
    0 & 1 & 1 \\
    0 & 0 & 1
\end{pmatrix}, \quad
c \mapsto \begin{pmatrix}
    1 & 0 & 1 \\
    0 & 1 & 0 \\
    0 & 0 & 1
\end{pmatrix}
\]

\subsection{Groups of order $16$}

Not lectured.

\begin{definition}
    Here are four new classes of groups of order $p^n$ for $n\geq 3$: \begin{itemize}
        \item $\Mod_{p^n}$
        \item $D_{2^n}$
        \item $SD_{2^n}$
        \item $Q_{2^n}$
    \end{itemize}
\end{definition}

\subsection{The transfer homomorphism}

\end{document}